% Volume IV: Optimization and Computation
% Continuity note for Chapter 21.
% Previous dependencies: Chapter 20 defines objectives, feasible sets, convexity, subgradients, regularization, conditioning, and local/global optimality; Chapter 7 defines shape-first gradients, Jacobians, Hessians, and matrix derivatives; Chapters 8-10 define Taylor approximation and local linear/quadratic models.
% New durable concepts: gradient descent, steepest descent, descent lemma, step-size stability, convergence on smooth convex objectives, condition-number effects, momentum, heavy-ball intuition, exponential moving averages, RMSProp and Adam-style adaptive scaling, Newton steps, damping, quasi-Newton secant intuition, mini-batch stochastic gradients, unbiasedness, variance, learning-rate schedules, and optimization noise.
% Forward hooks: Chapter 23 connects these updates to automatic differentiation; Chapter 24 studies numerical stability and large-scale linear algebra; deep-learning chapters reuse stochastic optimization, momentum, Adam, curvature, and training dynamics; neuroscience chapters reuse online learning and noisy adaptation ideas.
% Notation commitments: Parameters are theta in R^d; objectives are L(theta) or f(theta); gradients are g_k=nabla L(theta_k); Hessians are H_k=nabla^2 L(theta_k); step size is alpha or eta; momentum state is v_k; mini-batches are B_k; empirical risk is n^{-1} sum_i ell_i(theta).
% Figure plan: no figures are included in this authored source. Proposed raster figures are listed in imagegen-figures/ch21-figure-metadata.md.
\section{Chapter 21: Gradient Methods, Newton Methods, and Stochastic Optimization}
\label{ch:gradient-newton-stochastic-optimization}

Chapter 20 described the geometry of optimization problems. This chapter studies algorithms that move through that geometry. An optimizer is a dynamical system in parameter space: it observes local information such as gradients, curvature, or noisy mini-batch estimates, then chooses the next parameter vector.

The central tension is simple. Small steps are stable but slow. Large steps can be fast but may overshoot. Curvature information can dramatically improve steps but may be expensive. Stochastic gradients make large datasets trainable but inject noise. Modern machine learning and computational neuroscience use all of these tradeoffs.

\paragraph{Chapter problem.}
How do first-order, momentum, adaptive, second-order, and stochastic optimization algorithms use local information to reduce objectives, and what can go wrong when steps, curvature, or noise are mishandled?

\paragraph{Section map.}
Section 21.1 derives gradient descent from local linear approximation and studies step-size stability. Section 21.2 introduces momentum and adaptive methods as stateful modifications of gradient descent. Section 21.3 derives Newton and quasi-Newton methods from local quadratic approximation. Section 21.4 develops stochastic gradient descent from empirical-risk sums, mini-batch estimators, and gradient noise.

\subsection{21.1 Gradient descent}

\paragraph{Guiding question.}
Why does moving in the negative gradient direction reduce a differentiable objective, and how does the step size control stability?

\paragraph{Beginner-facing intuition.}
The gradient points in the direction of fastest local increase under the Euclidean norm. Therefore the negative gradient points in the direction of fastest local decrease. Gradient descent repeatedly asks: at the current point, which local direction lowers the objective most, and how far should we move?

The phrase ``local'' is essential. The gradient is computed at the current point. If the step is too large, the objective may curve before the next point is reached, and the local prediction can fail. Step-size choice is therefore part of the method, not an implementation detail.

\paragraph{Formal definitions and notation.}
Let \(L:\mathbb R^d\to\mathbb R\) be differentiable. Gradient descent defines an iterative sequence
\[
\boxed{
\theta_{k+1}=\theta_k-\alpha_k\nabla L(\theta_k),
}
\]
where \(\alpha_k>0\) is the step size or learning rate. If \(g_k=\nabla L(\theta_k)\), the update is \(\theta_{k+1}=\theta_k-\alpha_k g_k\).

The steepest-descent statement follows from the first-order approximation
\[
L(\theta+\Delta)\approx L(\theta)+\nabla L(\theta)^\top \Delta.
\]
Among all directions \(\Delta\) with \(\|\Delta\|_2=\varepsilon\), the inner product is minimized when
\[
\Delta=-\varepsilon\frac{\nabla L(\theta)}{\|\nabla L(\theta)\|_2},
\]
by Cauchy-Schwarz. This is the negative gradient direction.

\paragraph{Core mechanism: descent lemma and scalar stability.}
Assume \(L\) has \(M\)-Lipschitz gradient, meaning
\[
\|\nabla L(x)-\nabla L(y)\|_2\le M\|x-y\|_2.
\]
Then the descent lemma states
\[
L(y)\le L(x)+\nabla L(x)^\top(y-x)+\frac{M}{2}\|y-x\|_2^2.
\]
Substitute \(y=x-\alpha\nabla L(x)\):
\[
L(x-\alpha\nabla L(x))
\le
L(x)-\alpha\|\nabla L(x)\|_2^2
+\frac{M\alpha^2}{2}\|\nabla L(x)\|_2^2.
\]
Thus
\[
L(x-\alpha\nabla L(x))
\le
L(x)-\alpha\left(1-\frac{M\alpha}{2}\right)\|\nabla L(x)\|_2^2.
\]
If \(0<\alpha<2/M\), the right side decreases whenever the gradient is nonzero.

The scalar quadratic shows the same stability exactly. Let
\[
L(\theta)=\frac12 a\theta^2,\qquad a>0.
\]
Gradient descent gives
\[
\theta_{k+1}=(1-\alpha a)\theta_k.
\]
The sequence converges to \(0\) if and only if
\[
|1-\alpha a|<1,
\]
which is equivalent to
\[
\boxed{0<\alpha<\frac{2}{a}.}
\]
If \(\alpha\) is too large, iterates oscillate and diverge.

\paragraph{Concrete worked example.}
Let \(L(\theta)=(\theta-4)^2\). Then \(L'(\theta)=2(\theta-4)\). With \(\alpha=0.25\) and \(\theta_0=0\),
\[
\theta_1=0-0.25(-8)=2,
\]
\[
\theta_2=2-0.25(-4)=3,
\]
\[
\theta_3=3-0.25(-2)=3.5.
\]
The iterates approach \(4\). With \(\alpha=1.1\), the update factor for the quadratic curvature \(a=2\) is \(1-\alpha a=-1.2\), whose magnitude exceeds \(1\), so the iteration diverges by alternating around the minimum with increasing amplitude.

\paragraph{Computational neuroscience example.}
When fitting a firing-rate model by minimizing squared prediction error, the gradient with respect to a gain parameter tells how increasing that gain would change the loss on the current dataset. Gradient descent adjusts the gain in the direction that locally reduces mismatch. If the gain affects firing rates strongly, the curvature can be large, and a step size that worked for another parameter may be unstable.

\paragraph{AI and machine-learning example.}
In a neural network, automatic differentiation computes \(\nabla L(\theta)\) for millions of parameters. Gradient descent applies the same update rule coordinatewise. In practice, full-batch gradient descent is often replaced by mini-batch variants, but the local Taylor logic remains the same: gradients are first-order local decrease directions for the training objective.

\paragraph{Common misconceptions and failure modes.}
\begin{itemize}[leftmargin=*]
\item The negative gradient is steepest only relative to the chosen norm and parameterization.
\item A smaller learning rate is not always better; it may make progress impractically slow.
\item A decreasing objective does not guarantee good generalization.
\item A zero gradient can be a saddle point in nonconvex problems.
\item Rescaling parameters can change the behavior of plain gradient descent because the Euclidean geometry changes.
\end{itemize}

\paragraph{Exercises.}
\begin{enumerate}[label=\textbf{21.1.\arabic*.}, leftmargin=*]
\item \textbf{Mechanical.} For \(L(\theta)=3\theta^2\), write the gradient-descent update and find the range of constant \(\alpha\) for convergence.
\item \textbf{Proof or derivation.} Use Cauchy-Schwarz to show that, among directions of fixed Euclidean norm, \(-\nabla L(\theta)\) gives the largest first-order decrease.
\item \textbf{Numerical/computational.} Starting at \(\theta_0=10\), run three gradient-descent steps on \(L(\theta)=(\theta-1)^2\) with \(\alpha=0.1\).
\item \textbf{Interpretation.} In a neural model with parameters measured on different physical scales, why might one global step size be inefficient?
\item \textbf{Misconception/debugging.} A student says, ``If the objective went up after a gradient step, the gradient must have been wrong.'' Give another explanation.
\end{enumerate}

\subsection{21.2 Momentum and adaptive methods}

\paragraph{Guiding question.}
How do optimizers use memory of past gradients or per-coordinate scaling to move through ravines, noise, and uneven curvature?

\paragraph{Beginner-facing intuition.}
Plain gradient descent reacts only to the current gradient. Momentum adds a velocity state, so repeated gradients in the same direction accumulate and alternating gradients cancel. This helps in narrow valleys: the optimizer can move consistently along the valley while damping oscillations across it.

Adaptive methods change the scale of each coordinate's step using recent gradient magnitudes. A coordinate with consistently large gradients receives smaller effective steps; a coordinate with small gradients receives larger effective steps. This can help when parameters have different units, frequencies, or curvature scales.

\paragraph{Formal definitions and notation.}
One common momentum recursion is
\[
v_{k+1}=\beta v_k+\nabla L(\theta_k),
\qquad
\theta_{k+1}=\theta_k-\alpha v_{k+1},
\]
where \(0\le \beta<1\). The state \(v_k\) is an exponentially weighted sum of past gradients:
\[
v_{k+1}=g_k+\beta g_{k-1}+\beta^2 g_{k-2}+\cdots,
\]
ignoring initialization transients.

RMSProp-style scaling keeps a running second moment
\[
s_{k+1}=\rho s_k+(1-\rho)g_k\odot g_k,
\]
where \(\odot\) is elementwise multiplication, and updates
\[
\theta_{k+1}
=\theta_k-\alpha\frac{g_k}{\sqrt{s_{k+1}}+\varepsilon}
\]
elementwise.

Adam combines first- and second-moment estimates:
\[
m_{k+1}=\beta_1m_k+(1-\beta_1)g_k,
\]
\[
s_{k+1}=\beta_2s_k+(1-\beta_2)g_k\odot g_k,
\]
with bias-corrected versions often used in practice before the elementwise scaled update. The small \(\varepsilon>0\) prevents division by zero.

\paragraph{Core mechanism: ravine-shaped quadratic.}
Consider
\[
L(x,y)=\frac12(100x^2+y^2).
\]
The gradient is
\[
\nabla L(x,y)=(100x,y).
\]
The \(x\)-direction is steep, while the \(y\)-direction is shallow. Plain gradient descent with a stable step size must respect the steep direction:
\[
\alpha<\frac{2}{100}=0.02.
\]
That makes progress along \(y\) slow because the shallow curvature is \(1\). Momentum can accumulate consistent motion along \(y\), while the sign-changing \(x\)-updates partly cancel. Adaptive scaling divides by recent gradient magnitudes, reducing the effective step in the steep coordinate.

These mechanisms are useful but not guarantees. Momentum can overshoot if \(\alpha\) and \(\beta\) are too large. Adaptive methods can behave differently from optimizing the original objective under a fixed geometry because they change the coordinate-wise metric over time.

\paragraph{Concrete worked example.}
Let \(L(\theta)=\frac12\theta^2\), \(\theta_0=4\), \(v_0=0\), \(\alpha=0.1\), and \(\beta=0.9\). Then \(g_0=4\),
\[
v_1=0.9(0)+4=4,\qquad \theta_1=4-0.1(4)=3.6.
\]
Next \(g_1=3.6\),
\[
v_2=0.9(4)+3.6=7.2,\qquad \theta_2=3.6-0.1(7.2)=2.88.
\]
Momentum moves farther than plain gradient descent at the second step because gradients have pointed in the same direction.

\paragraph{Computational neuroscience example.}
When fitting a recurrent neural dynamics model, gradients for some parameters can oscillate because changing a recurrent weight affects future states repeatedly. Momentum can smooth these oscillations in parameter space. This is an optimization device; it should not be interpreted as a literal claim that biological synapses implement the same velocity recursion.

\paragraph{AI and machine-learning example.}
Adam and related methods are common default optimizers for deep networks because they handle noisy mini-batch gradients and coordinate scale differences reasonably well. However, weight decay, learning-rate schedules, batch size, and normalization layers interact with adaptive scaling. A model that trains quickly with Adam may still need validation checks, calibration checks, or later fine-tuning with another optimizer.

\paragraph{Common misconceptions and failure modes.}
\begin{itemize}[leftmargin=*]
\item Momentum is not simply a larger learning rate; it is a stateful recursion with memory.
\item Adaptive scaling does not remove the need to tune the base learning rate.
\item Adam's fast training loss decrease does not guarantee best test performance.
\item Per-coordinate scaling is parameterization-dependent.
\item Large momentum can cause oscillation or delayed reaction when the gradient direction changes.
\end{itemize}

\paragraph{Exercises.}
\begin{enumerate}[label=\textbf{21.2.\arabic*.}, leftmargin=*]
\item \textbf{Mechanical.} With \(g_0=2\), \(g_1=1\), \(v_0=0\), and \(\beta=0.5\), compute \(v_1\) and \(v_2\) for the momentum recursion \(v_{k+1}=\beta v_k+g_k\).
\item \textbf{Proof or derivation.} Expand \(v_{k+1}=\beta v_k+g_k\) recursively to show that it is an exponentially weighted sum of past gradients.
\item \textbf{Numerical/computational.} For \(L(x,y)=\frac12(100x^2+y^2)\), compute the gradient at \((1,1)\). Compare the plain gradient step with \(\alpha=0.01\) to an RMSProp-style step with \(s=(10000,1)\), \(\varepsilon=0\), and \(\alpha=0.1\).
\item \textbf{Interpretation.} Why can momentum help in a long narrow valley but hurt when the optimum is near and the velocity remains large?
\item \textbf{Misconception/debugging.} A student says, ``Adam is scale invariant, so initialization and normalization no longer matter.'' Explain why this overstates the method.
\end{enumerate}

\subsection{21.3 Newton and quasi-Newton methods}

\paragraph{Guiding question.}
How can second-order curvature information choose better-scaled steps than the gradient alone?

\paragraph{Beginner-facing intuition.}
The gradient tells which way is downhill locally, but it does not tell how curved the objective is. Newton's method uses a local quadratic model. If the objective is steep in one direction and shallow in another, Newton's step rescales the gradient by curvature, taking smaller steps in steep directions and larger steps in shallow directions.

The cost is that curvature is expensive. A Hessian for \(d\) parameters has \(d^2\) entries, and solving a Hessian system can be costly. Quasi-Newton methods build approximate curvature from changes in gradients rather than forming the true Hessian.

\paragraph{Formal definitions and notation.}
Let \(L:\mathbb R^d\to\mathbb R\) be twice differentiable. Around \(\theta_k\), the second-order Taylor model is
\[
m_k(p)
=L(\theta_k)+g_k^\top p+\frac12p^\top H_kp,
\]
where
\[
g_k=\nabla L(\theta_k),\qquad H_k=\nabla^2L(\theta_k).
\]
If \(H_k\) is positive definite, the minimizer of \(m_k\) satisfies
\[
\nabla_p m_k(p)=g_k+H_kp=0.
\]
Thus the Newton step \(p_k\) solves
\[
\boxed{
H_kp_k=-g_k,
}
\]
and the update is
\[
\theta_{k+1}=\theta_k+p_k
=\theta_k-H_k^{-1}g_k.
\]
In practice one solves the linear system rather than forming \(H_k^{-1}\).

Damped Newton uses
\[
\theta_{k+1}=\theta_k+\alpha_k p_k
\]
with a line search or trust region. If \(H_k\) is indefinite, the pure Newton direction may not be a descent direction.

\paragraph{Core mechanism: exact solution for quadratics and quasi-Newton secants.}
For a positive definite quadratic
\[
L(\theta)=\frac12\theta^\top A\theta-b^\top\theta+c,
\]
the gradient and Hessian are
\[
g(\theta)=A\theta-b,\qquad H(\theta)=A.
\]
Newton's step solves
\[
Ap=-(A\theta_k-b),
\]
so
\[
\theta_{k+1}=\theta_k+p=A^{-1}b,
\]
the exact minimizer, in one step. This explains the appeal of curvature: it solves the local quadratic model exactly.

Quasi-Newton methods approximate the Hessian or its inverse. Let
\[
s_k=\theta_{k+1}-\theta_k,\qquad
y_k=g_{k+1}-g_k.
\]
For a true constant Hessian, \(y_k=Hs_k\). Quasi-Newton updates choose a matrix \(B_{k+1}\) satisfying the secant condition
\[
B_{k+1}s_k=y_k
\]
while staying close to the previous approximation \(B_k\). BFGS and L-BFGS are widely used examples. L-BFGS stores only a limited history of \(s_k,y_k\) pairs, making it practical for larger problems.

\paragraph{Concrete worked example.}
Let
\[
L(x,y)=\frac12(10x^2+y^2).
\]
At \((x,y)=(1,1)\),
\[
g=(10,1)^\top,
\qquad
H=
\begin{bmatrix}
10&0\\
0&1
\end{bmatrix}.
\]
The Newton step solves
\[
Hp=-g,
\]
so
\[
p=(-1,-1)^\top.
\]
The next point is \((0,0)\), the minimizer. A gradient step with \(\alpha=0.1\) gives
\[
(1,1)-0.1(10,1)=(0,0.9),
\]
which fixes the steep coordinate but moves slowly in the shallow coordinate. Newton rescales the two coordinates by curvature.

\paragraph{Computational neuroscience example.}
In a Poisson generalized linear model for spike counts, the negative log-likelihood has gradients and Hessians that depend on stimulus filters and coupling weights. Newton or quasi-Newton methods can exploit curvature to fit moderate-sized models quickly. For very large neural populations or long time histories, exact Hessians may be too expensive, so L-BFGS, Hessian-vector products, or first-order methods are used.

\paragraph{AI and machine-learning example.}
Deep networks have enormous parameter vectors, so full Newton steps are usually impractical. However, second-order ideas appear in natural gradient, K-FAC-style approximations, Gauss-Newton methods, L-BFGS for smaller models, and curvature-aware learning-rate schedules. The goal is the same: use curvature information to improve scaling and avoid inefficient zig-zagging.

\paragraph{Common misconceptions and failure modes.}
\begin{itemize}[leftmargin=*]
\item Newton's method is not automatically safer than gradient descent; an indefinite Hessian can point uphill.
\item Computing \(H^{-1}\) explicitly is usually unnecessary and numerically poor; solve \(Hp=-g\).
\item Quadratic convergence is a local result requiring regularity and a good initial point.
\item Quasi-Newton methods approximate curvature; they are not exact Newton methods.
\item A Hessian with small eigenvalues signals flat or weakly identified directions where steps and uncertainty can be large.
\end{itemize}

\paragraph{Exercises.}
\begin{enumerate}[label=\textbf{21.3.\arabic*.}, leftmargin=*]
\item \textbf{Mechanical.} For \(L(\theta)=\frac12 a\theta^2-b\theta\) with \(a>0\), write the Newton update and show it reaches \(b/a\) in one step.
\item \textbf{Proof or derivation.} Derive the Newton step by minimizing the second-order Taylor model \(m_k(p)\).
\item \textbf{Numerical/computational.} At \(\theta=(1,2)\), let \(g=(4,2)\) and \(H=\begin{bmatrix}4&0\\0&1\end{bmatrix}\). Compute the Newton step and the updated point.
\item \textbf{Interpretation.} Why might a Hessian be poorly conditioned when fitting a neural model with highly correlated stimulus features?
\item \textbf{Misconception/debugging.} A student writes \(p=-H^{-1}g\) and computes an explicit inverse in code. Explain the preferred numerical approach and why.
\end{enumerate}

\subsection{21.4 Stochastic gradient descent}

\paragraph{Guiding question.}
How can we optimize an average loss using random mini-batches, and what does the resulting gradient noise do?

\paragraph{Beginner-facing intuition.}
Many objectives are averages over large datasets:
\[
L(\theta)=\frac1n\sum_{i=1}^n \ell_i(\theta).
\]
Computing the full gradient every step can be expensive. Stochastic gradient descent estimates the gradient from one example or a mini-batch. The estimate is noisy, but much cheaper. The algorithm trades exact local information for many inexpensive updates.

The noise is not merely a nuisance. It can help escape shallow regions, interact with learning-rate schedules, and affect generalization. But too much noise or too large a learning rate can prevent convergence.

\paragraph{Formal definitions and notation.}
Let
\[
L(\theta)=\frac1n\sum_{i=1}^n \ell_i(\theta),
\]
where each \(\ell_i\) is the loss on example \(i\). The full gradient is
\[
\nabla L(\theta)=\frac1n\sum_{i=1}^n \nabla \ell_i(\theta).
\]
Let \(B_k\subset\{1,\ldots,n\}\) be a mini-batch sampled uniformly, with \(|B_k|=b\). The mini-batch gradient estimator is
\[
\widehat g_k
=\frac1b\sum_{i\in B_k}\nabla \ell_i(\theta_k).
\]
When sampled uniformly, with appropriate handling of replacement or without-replacement details,
\[
\boxed{
\mathbb E[\widehat g_k\mid \theta_k]=\nabla L(\theta_k).
}
\]
The SGD update is
\[
\boxed{
\theta_{k+1}=\theta_k-\alpha_k\widehat g_k.
}
\]

\paragraph{Core mechanism: unbiasedness and variance.}
If mini-batch indices \(I_1,\ldots,I_b\) are sampled independently and uniformly from \(\{1,\ldots,n\}\), define
\[
G_j=\nabla \ell_{I_j}(\theta).
\]
Then
\[
\mathbb E[G_j]=\frac1n\sum_{i=1}^n\nabla \ell_i(\theta)=\nabla L(\theta).
\]
Therefore
\[
\mathbb E[\widehat g]=
\mathbb E\left[\frac1b\sum_{j=1}^bG_j\right]
=\nabla L(\theta).
\]
If the covariance of a single-example gradient is \(\Sigma_g\), independent sampling gives
\[
\operatorname{Cov}(\widehat g)=\frac{1}{b}\Sigma_g.
\]
Larger mini-batches reduce gradient variance, but they cost more per update and may change the training dynamics.

Classical stochastic approximation uses decreasing step sizes satisfying
\[
\sum_{k=0}^\infty \alpha_k=\infty,
\qquad
\sum_{k=0}^\infty \alpha_k^2<\infty
\]
under appropriate assumptions. Deep learning often uses piecewise-constant or scheduled learning rates instead, chosen empirically.

\paragraph{Concrete worked example.}
Let
\[
\ell_1(\theta)=\frac12(\theta-1)^2,\qquad
\ell_2(\theta)=\frac12(\theta-3)^2.
\]
Then
\[
L(\theta)=\frac12\ell_1(\theta)+\frac12\ell_2(\theta),
\]
and
\[
\nabla L(\theta)=\frac12(\theta-1)+\frac12(\theta-3)=\theta-2.
\]
At \(\theta_0=0\), the full gradient is \(-2\). A full gradient step with \(\alpha=0.1\) gives \(\theta_1=0.2\).

If SGD samples only example 1, the gradient is \(-1\), and the update gives \(\theta_1=0.1\). If it samples example 2, the gradient is \(-3\), and the update gives \(\theta_1=0.3\). The average of these two possible SGD updates is \(0.2\), matching the full-gradient update in expectation.

\paragraph{Computational neuroscience example.}
Online fitting of a neural encoding model can update parameters after each trial or small block of trials. The gradient from one trial is noisy because spikes are variable, but it provides immediate information. This resembles stochastic approximation: synaptic or model parameters are adjusted using noisy samples of a larger expected objective. The mathematical analogy is useful, but it does not imply that biological plasticity literally optimizes the exact loss used by the model.

\paragraph{AI and machine-learning example.}
Training deep networks almost always uses mini-batch stochastic gradients. The mini-batch size influences hardware efficiency, gradient variance, and the effective noise scale. Learning-rate warmup, decay schedules, momentum, Adam, and weight decay are all built around managing this noisy iterative process.

\paragraph{Common misconceptions and failure modes.}
\begin{itemize}[leftmargin=*]
\item A stochastic gradient is not the true gradient; it is an estimator whose noise depends on the mini-batch.
\item Unbiased does not mean low variance.
\item Larger batches reduce gradient noise but do not automatically improve generalization or wall-clock training.
\item Shuffling and data dependence matter; time-correlated neural data may violate iid mini-batch assumptions.
\item Training loss curves under SGD are noisy, so single-step increases are not necessarily failures.
\end{itemize}

\paragraph{Exercises.}
\begin{enumerate}[label=\textbf{21.4.\arabic*.}, leftmargin=*]
\item \textbf{Mechanical.} For \(L(\theta)=n^{-1}\sum_{i=1}^n\ell_i(\theta)\), write the full gradient and a mini-batch gradient estimator of batch size \(b\).
\item \textbf{Proof or derivation.} Prove that the mini-batch gradient estimator is unbiased when indices are sampled uniformly with replacement.
\item \textbf{Numerical/computational.} In the two-example worked example, compute the SGD update from \(\theta_0=1\) with \(\alpha=0.2\) if example 1 is sampled, and if example 2 is sampled.
\item \textbf{Interpretation.} Why might mini-batch gradients from consecutive trials in a neuroscience experiment be more correlated than mini-batches sampled from shuffled images?
\item \textbf{Misconception/debugging.} A student says, ``Because SGD is noisy, it cannot converge.'' Explain how decreasing or scheduled learning rates and variance reduction by batching address this.
\end{enumerate}

\paragraph{Closing bridge.}
Gradient, momentum, adaptive, Newton, and stochastic methods are all ways to convert local information into parameter motion. Chapter 22 adds another layer: many objectives include constraints, dual variables, or entire functions as decision objects. Those settings require Lagrangians, KKT conditions, duality, and variational calculus.

\clearpage
